% !TeX root = main.tex
\chapter*{Preface}\addcontentsline{toc}{chapter}{Preface}\markboth{Preface}{}
A solution of the wave equation is sometimes written as a smooth function, sometimes as a point in a phase space, sometimes as a label for a quantum observable. These descriptions cannot simply be substituted for one another. A weak solution may not even have a value at a particular point. A label for an observable can be changed by the equation of motion without changing the observable. A commutator does not assign that observable a vacuum expectation value. How does one pass from each object to the next, and what survives the passage?

This book develops the mathematics of those passages. It can be read independently of \emph{Lorentzian Sewing: Geometry, Dynamics, and Observables}. The reader is assumed to know sets and functions, elementary linear algebra over the real and complex numbers, and calculus in Euclidean space, including change of variables and integration by parts. Familiarity with waves and an action principle will make some applications recognizable, but neither a course in differential geometry nor one in functional analysis is assumed. The inverse and implicit function theorems in Euclidean space, local existence and smooth dependence for smooth ordinary differential equations, and the elementary results of Lebesgue integration are available as analysis inputs. Where we require a deeper theorem we state the version being used and identify its source. A proof in a flat example is not a proof of that theorem on an arbitrary curved space.

Each change of mathematical object answers a question. Charts give local formulas for a form; integration must return a number independent of those charts. Field values live in different fibers; a connection makes their derivatives meaningful. When smoothness is weakened, even restriction to a boundary becomes a theorem with a sharp threshold. Passing from the classical bracket to a quantum algebra leaves two further questions open: which states are positive, and which singular fields may be multiplied?

Each part develops its own arguments before applying them to interfaces. The exercises ask, among other things, for counterexamples: a definition is often understood for the first time upon seeing which of its hypotheses one cannot drop. The last chapter gives a second way into the book for a reader with experience: work through one field problem, and follow the references back at the first operation one cannot justify. A guide at the back locates the mathematics needed for each chapter of \emph{Lorentzian Sewing}.

\paragraph{Conventions.} A dimension counts spacetime coordinates unless another space is named. The metric signature is $(-,+,\ldots,+)$. The letter $\dd$ differentiates on a base manifold; $\delta$ differentiates with respect to fields. A compactly supported field variation need not vanish at an artificial internal cut. We will distinguish ordinary exterior boundaries, where a variational problem has to be specified, from a cut introduced only to describe the same object using smaller regions.
